\section*{Chapter \ref{ch:rc:regulatory-capital-for-trading-books} --- Regulatory Capital for Trading Books}

\begin{solution}{exo:rc:regulatory-capital-for-trading-books:1}
$\max(e^{-0.03\times8/2},0.40) = e^{-0.12} = 88.7\%$; high $\min(1.25\times0.887,1) = 100\%$; low
$\max(2\times0.887-1,0.75\times0.887) = 77.4\%$.
\end{solution}

\begin{solution}{exo:rc:regulatory-capital-for-trading-books:2}
The book is long the ten-year and short the two-year: the positions offset through their positive
correlation, and less correlation means less offset. The same happens across the long-euro and
short-yen FX buckets. With the vega and curvature charges unchanged, the low scenario adds most.
\end{solution}

\begin{solution}{exo:rc:regulatory-capital-for-trading-books:3}
Amber: the correlation passes (above 0.80) but the distribution distance (0.10) lies between 0.09 and
0.12. The desk keeps the internal model but pays the surcharge.
\end{solution}

\begin{solution}{exo:rc:regulatory-capital-for-trading-books:4}
The standardised approach charges the straddle's curvature (a $\pm10.6\%$ EURUSD shock) and its vega at a
100\% weight, sizes calibrated for any bank's book; the internal model measures the ten-day ES of the
actual book, where EURUSD moves are far smaller, even in the stressed year, and diversification
applies.
\end{solution}

\begin{solution}{exo:rc:regulatory-capital-for-trading-books:5}
No: it has more than 24 observations, but the last eight months include 90-day periods with none, so
criterion (1) fails, and 30 is below the 100 of criterion (2). It is non-modellable.
\end{solution}

\begin{solution}{exo:rc:regulatory-capital-for-trading-books:6}
The diversified ES relies on the model's cross-class correlations, which can break in stress; the sum
over classes assumes none. Weighting them half and half limits the capital benefit of
cross-class diversification.
\end{solution}

\begin{solution}{exo:rc:regulatory-capital-for-trading-books:7}
$\sqrt{5^2+\bigl(3\sqrt{(20-10)/10}\bigr)^2} = \sqrt{34} = 5.83$.
\end{solution}

\begin{solution}{exo:rc:regulatory-capital-for-trading-books:8}
Lower is not better: the model must pass backtesting and attribution desk by desk, its stressed
calibration and non-modellable factors add capital, and the standardised approach remains a
benchmark (and in some jurisdictions a floor). A low number with a failing attribution test is
evidence against the model.
\end{solution}

\begin{solution}{pb:rc:regulatory-capital-for-trading-books:1}
\textbf{1.} USD~68.29 million, in the low correlation scenario.
\textbf{2.} USD~25.56 million: 1.5 times an IMCC of 17.04 million, from a stressed ES of 14.41 million
for the book and 5.67 and 14.01 million for rates and FX.
\textbf{3.} FX curvature, USD~39.23 million, from the short straddle.
\textbf{4.} It contains the dollar's surge and the euro's fall below parity, the largest EURUSD moves in
the sample, which the short straddle and the euro position feel most.
\textbf{5.} Nothing in the internal model here, since volatility is not a modelled factor; in a model with
it, a 40-day horizon would scale up the vega-driven ES.
\textbf{6.} 0.703 and 0.044.
\textbf{7.} Amber: the correlation is between 0.70 and 0.80; the distance passes.
\textbf{8.} A surcharge of $0.5\times(68.29-25.56) = \text{USD}~21.36$ million; capital USD~46.93 million.
\textbf{9.} The standardised charge, USD~68.29 million.
\textbf{10.} The delta--gamma and no-yen models pass (green); the no-two-year model is amber through its
distance (0.100).
\textbf{11.} The proxy moves the two tenors together, so it misses the days the curve twists: the order of
the daily P\&Ls (ranks) is scrambled, while the spread of the distribution stays similar.
\textbf{12.} Real price observations of ten-year swaps or bonds: at least 24 a year with no 90-day gap
holding fewer than four.
\textbf{13.} Partly: a regression scales the move but still cannot see the twist; only a separate factor
restores the ranks.
\textbf{14.} Time to gather data and fix the model before the test affects capital.
\textbf{15.} Compare the fix's cost with the surcharge's present value over the desk's horizon, and the
risk of falling to red.
\textbf{16.} Backtesting checks the tail count; attribution checks that the risk model sees the same
P\&L drivers every day. A model can pass one and fail the other.
\textbf{17.} To improve risk models, or to game the statistics by choosing proxies that preserve ranks.
\textbf{18.} Where the rules and multipliers differ, a bank may book trades in the jurisdiction with the
lighter treatment, or run different models.
\textbf{19.} Named result: \emph{the desk that failed attribution}: internal-model capital USD~25.56 million,
standardised USD~68.29 million; Spearman 0.703 and KS 0.044 put it in amber, and capital becomes USD~46.93
million.
\textbf{20.} That the risk model and the pricing systems see the same world.
\end{solution}

\begin{solution}{iq:rc:regulatory-capital-for-trading-books:1}
VaR and stressed VaR at 99\% replaced by \omterm{def:rc:market-risk-measures:es}{expected shortfall} at 97.5\% with \omterm{def:rc:regulatory-capital-for-trading-books:ima}{liquidity horizons} and stressed
calibration; a sensitivity-based standardised approach; desk-level approval with backtesting and
P\&L attribution; \omterm{def:rc:regulatory-capital-for-trading-books:nmrf}{non-modellable risk factors}; a stricter book boundary.
\iqlookfor{ES, \omterm{def:rc:regulatory-capital-for-trading-books:ima}{liquidity horizons}, SBM and desk-level tests.}
\end{solution}

\begin{solution}{iq:rc:regulatory-capital-for-trading-books:2}
PV01 by currency and tenor; risk weights per tenor (divided by $\sqrt2$ for major currencies);
aggregation within the currency with tenor correlations and across currencies at 50\%; the three
correlation scenarios; curvature from parallel shocks if there are options.
\iqlookfor{the steps and parameters.}
\end{solution}

\begin{solution}{iq:rc:regulatory-capital-for-trading-books:3}
It compares the risk model's P\&L with the front office's by rank correlation and a distribution
distance. Desks fail from missing or proxied \omterm{def:rc:market-risk-measures:factor}{risk factors}, different data sources or snapshot times,
simplified revaluation, and basis risks.
\iqlookfor{the metrics and the common causes.}
\end{solution}

\begin{solution}{iq:rc:regulatory-capital-for-trading-books:4}
ES sees the size of tail losses and is coherent; \omterm{def:rc:regulatory-capital-for-trading-books:ima}{liquidity horizons} reflect that some positions take
weeks to exit, which a uniform ten-day VaR ignores; stressed calibration avoids capital falling in calm
markets.
\iqlookfor{tail, liquidity and procyclicality.}
\end{solution}

\begin{solution}{iq:rc:regulatory-capital-for-trading-books:5}
By the capital saving against the cost of models, data and controls, and by the likelihood of passing
backtesting, attribution and eligibility tests; desks with liquid, well-modelled risks first.
\iqlookfor{cost-benefit and test risk.}
\end{solution}

\begin{solution}{iq:rc:regulatory-capital-for-trading-books:6}
Long histories of real prices per \omterm{def:rc:market-risk-measures:factor}{risk factor} with observation counts; consistent market data for
front office and risk; daily HPL and RTPL by desk; full revaluation infrastructure for ES over many
scenarios and liquidity subsets; stress-period search; reporting and audit trails.
\iqlookfor{data, revaluation and controls.}
\end{solution}
